% Глава: 08-packages
% Статус: наполнено (первый прогон).
\chapter{Пакетная система: dpkg и apt}
\label{ch:08-packages}

\section{Зачем нужна пакетная система}

Устанавливать программы «вручную» (скачать архив, распаковать,
поставить зависимости) неудобно и ненадёжно.
\textbf{Пакетная система} хранит программы в стандартизированном
виде и отслеживает: что установлено, какие файлы к чему
относятся, какие зависимости есть. В Debian/Ubuntu
основной формат --- \textbf{\fpath{.deb}}.

Устанавливаем пакеты в свою ВМ на сервере \fpath{pglab.nsu.ru}:
образ Ubuntu идёт «голый», и каждый добирает нужное сам.

\begin{figure}[htb]
  \centering
  \begin{tikzpicture}[node distance=1.4cm]
    \node[process] (repo) {Депозитарий\\[0.2em] (apt)};
    \node[data, below=of repo] (cache) {Кэш пакетов\\[0.2em] /var/cache/apt};
    \node[process, below=of cache] (dpkg) {dpkg\\[0.2em] (распаковка .deb)};
    \node[storage, below=of dpkg] (fs) {Файлы системы\\[0.2em] /usr, /etc};
    \draw[flow, <->] (repo) -- node[right]{установка} (cache);
    \draw[flow, <->] (cache) -- (dpkg);
    \draw[flow, <->] (dpkg) -- (fs);
  \end{tikzpicture}
  \caption{Связь между apt (репозиторий) и dpkg (развёртка).}
  \label{fig:pkg}
\end{figure}

\section{dpkg и apt: чем отличаются}

\begin{center}
\begin{tabular}{lp{0.4\linewidth}lp{0.4\linewidth}}
\toprule
\textbf{dpkg} & работает с конкретным файлом \fpath{.deb} и файлами
системы; зависимости сам не разрешает \\
\textbf{apt} & управляет набором пакетов с учётом зависимостей и репозиториев \\
\midrule
\multicolumn{2}{l}{На практике: \cmd{apt} --- интерфейс для пользователя,} \\
\multicolumn{2}{l}{\cmd{dpkg} --- низкоуровневый исполнитель.} \\
\bottomrule
\end{tabular}
\end{center}

Обычно работают через \cmd{apt} (или \cmd{apt-get}); \cmd{dpkg}
используют для точечных задач (установка локального \fpath{.deb},
просмотр файлов пакета).

\section{Основные операции}

\begin{minted}[bgcolor=gray!6]{bash}
sudo apt update              # обновить список доступных пакетов
sudo apt upgrade             # обновить установленные пакеты
sudo apt install <pkg>       # установить пакет
sudo apt remove <pkg>        # удалить, файлы конфигурации остаются
sudo apt purge <pkg>         # удалить вместе с конфигурацией
sudo apt autoremove          # удалить ненужные зависимости
apt search <term>            # поиск по каталогу
apt show <pkg>               # сведения о пакете
apt list --installed         # список установленных пакетов
\end{minted}

\textbf{update} и \textbf{upgrade} --- разные вещи:
\cmd{update} перечитывает списки доступных версий из
репозиториев; \cmd{upgrade} ставит новые версии уже
установленного.

\section{Работа с локальными \fpath{.deb}}

Если вам дали файл \fpath{package.deb} --- его можно установить
двумя способами:
\begin{minted}[bgcolor=gray!6]{bash}
sudo dpkg -i package.deb          # напрямую
sudo apt install ./package.deb    # apt разрешает зависимости
\end{minted}

Первый способ быстрее, но если у пакета есть
зависимости, которых нет в системе, \cmd{dpkg}
остановится. Тогда \cmd{apt install ./package.deb}
доустановит всё из репозиториев.

\section{Инспекция пакетов без установки}

\begin{minted}[bgcolor=gray!6]{bash}
dpkg-deb --info package.deb      # имя, версия, описание
dpkg-deb --contents package.deb  # список файлов внутри
dpkg-deb --show --showformat='...\\n' package.deb  # свои поля
\end{minted}

После установки:
\begin{minted}[bgcolor=gray!6]{bash}
dpkg -l                          # все пакеты и их состояние
dpkg -L <pkg>                    # файлы, установленные пакетом
dpkg -S <path>                   # какой пакет установил файл
dpkg -s <pkg>                    # подробная информация
dpkg -C                          # проверка состояния, поиск ошибок
\end{minted}

\section{Стандартные каталоги}

Пакеты в Debian/Ubuntu обычно ставят:
\begin{itemize}
  \item исполняемые файлы --- в \fpath{/usr/bin} и \fpath{/usr/sbin};
  \item библиотеки --- в \fpath{/usr/lib};
  \item документацию --- в \fpath{/usr/share/doc/<pkg>};
  \item конфигурацию --- в \fpath{/etc/<pkg>};
  \item метаданные --- в \fpath{/var/lib/dpkg/info}.
\end{itemize}

\cmd{dpkg -L} показывает реальный список для конкретного
пакета --- не полагайтесь на память.

\section{Типичные ошибки}
\begin{itemize}
  \item \emph{Unable to locate package} --- сначала
    \cmd{sudo apt update};
  \item Пакет установлен, но команда не находится ---
    посмотрите \cmd{dpkg -L}, какой бинарник поставил пакет;
  \item После удаления пакета остались файлы конфигурации ---
    это нормально; уберите \cmd{sudo apt purge <pkg>};
  \item Зависимости не разрешаются --- не используйте
    \cmd{dpkg -i} для пакетов с зависимостями;
  \item \cmd{apt} ругается на «held» пакеты --- посмотрите
    \cmd{apt-mark showhold}.
\end{itemize}

\section{Связь с остальным курсом}

Пакетная система нужна и для установки
\cmd{gcc}, \cmd{make}, \cmd{valgrind}
(\nameref{ch:13-c-compilation}, \nameref{ch:17-debugging}),
и для \cmd{git} (\nameref{ch:25-git}), и для \cmd{docker}
(\nameref{ch:19-docker-basics}). В Docker-образе пакетная
система --- тоже \cmd{apt} (если образ на Debian/Ubuntu),
см.\ \nameref{ch:20-dockerfile}.

\section{Задачи для самопроверки}
\begin{enumerate}
  \item Объясните разницу между \cmd{apt update} и
    \cmd{apt upgrade}. Что делает каждое?
  \item Чем \cmd{dpkg} отличается от \cmd{apt} в
    разрешении зависимостей?
  \item Как узнать, какие файлы принадлежат
    установленному пакету?
  \item Чем \cmd{remove} отличается от \cmd{purge}?
  \item Как установить локальный \fpath{.deb} так,
    чтобы зависимости разрешались автоматически?
\end{enumerate}
